\clearpage
\subsection{Beschränktheit: Werte bleiben zwischen festen Grenzen}
\label{subsec:real_functions_boundedness}
Gebrochen-rationale Funktionen müssen keine Polstellen haben.
Der Term $f(x)=1/(1+x^2)$ ist für alle reellen Eingaben definiert;
seine Werte bleiben in einem festen Bereich.
\begin{center}
  \CachedTikzFigure{content/resources/chapters/real_functions/tikz/functions_bounded.tex}
\end{center}
\par\noindent\begin{minipage}{\linewidth}
\begin{FunctionPropertyBox}{Beschränktheit}{real_functions_bounded}
  Eine Funktion $f:D\to\R$ heißt \emph{nach oben beschränkt}
  \EnglishTerm{bounded above}, wenn eine Zahl $M\in\R$ mit
  $f(x)\le M$ für alle $x\in D$ existiert.
  Sie heißt \emph{nach unten beschränkt} \EnglishTerm{bounded below},
  wenn entsprechend $m\le f(x)$ gilt.

  Sind beide Bedingungen erfüllt, heißt $f$ \emph{beschränkt}
  \EnglishTerm{bounded}. Gleichwertig existiert $K\ge0$ mit
  $|f(x)|\le K$ für alle $x\in D$.
\end{FunctionPropertyBox}
\begin{ExampleBox}[before skip=0.3\baselineskip]{Schranke oder angenommener Wert?}{real_functions_bounds_example}
  Bestimmen Sie die \textcolor{green!55!black}{Wertemenge} von $g(x)=1+4/(1+x^2)$ auf $\R$.
  \begin{ClozeArea}[34]{7}
    Aus $x^2\ge0$ und dem positiven Nenner folgt
    \[
      0<\frac1{1+x^2}\le1
      \quad\Longrightarrow\quad
      1<g(x)\le5.
    \]
    Der Wert $5$ wird bei $x=0$ angenommen. Der Wert $1$ wird nur
    im Grenzfall $x\to\pm\infty$ erreicht; er gehört nicht zur \textcolor{green!55!black}{Wertemenge}.
    Für jedes $y\in(1,5]$ liefert
    $x=\sqrt{4/(y-1)-1}$ eine zulässige Eingabe.
    Also ist $g(\R)=(1,5]$.
  \end{ClozeArea}
\end{ExampleBox}
\end{minipage}
% Handschrift: PDF-Seite 30 (obere und untere Schranke im Originalbeispiel).
\begin{center}
  \CachedTikzFigure{content/resources/chapters/real_functions/tikz/functions_bounds_example.tex}
\end{center}
Für $g(x)=1+4/(1+x^2)$ sind also $m=1$ eine untere und $M=5$ eine
obere Schranke: $1<g(x)\le5$ für alle $x\in\R$. Eine Schranke muss
nicht selbst angenommen werden; die untere Schranke $1$ wird nur angenähert.

\begin{KeyPointBox}{Beschränkt heißt nicht: besitzt einen Grenzwert}{real_functions_bounds_not_limit}
  Beschränktheit begrenzt die möglichen Werte.
  Sie sagt noch nicht, ob sich diese Werte bei einer Annäherung auf
  einen einzigen Wert einstellen. Bei periodischen Funktionen werden
  wir diesen Unterschied wiedersehen.
\end{KeyPointBox}
